\section{Conclusion}

\label{sec:conclusion}

Active Semantic Optimization demonstrates that training-free adaptation via
TF-GRPO and Product-of-Experts decoding can achieve competitive performance on
challenging code generation benchmarks at a fraction of the cost of fine-tuning.
The key contributions are:

\begin{enumerate}[leftmargin=1.5em]
    \item \textbf{TF-GRPO:} A training-free extension of GRPO that extracts
          semantic advantages from grouped rollouts without gradient computation.
    \item \textbf{PoE Decoding:} A Bayesian framework for incorporating
          experiential priors as multiplicative token-level factors with
          theoretically optimal expert weights.
    \item \textbf{1-Bit Compression:} A quantization scheme achieving
          $29.5\times$ storage savings with bounded approximation error.
    \item \textbf{Convergence Guarantees:} Formal proofs that ASO converges to
          the Bayes-optimal predictor with explicit finite-sample bounds.
    \item \textbf{Cost Efficiency:} \$18 adaptation cost compared to
          \$10{,}000+ for parameter-efficient fine-tuning.
\end{enumerate}

\paragraph{Limitations.}
ASO requires a reward signal (test suite) for TF-GRPO, which is not available
for all tasks. The PoE framework assumes token-level decomposability, which may
not capture long-range dependencies. Performance still trails full fine-tuning
by 4--6 percentage points on SWE-bench.

\paragraph{Future Work.}
Multi-agent prior sharing via the DSO protocol (see companion paper),
extension to non-code domains (data science, DevOps, security), and
investigation of active learning strategies for prior acquisition.

% ============================================================================
% APPENDIX
% ============================================================================
\appendix

